\Titular%
{historia}%
{Buracos negros antes de Einstein?}%
{Manuel Vázquez Carreira}%
{Como Michell e Laplace teorizaron por primeira vez os buracos negros.}%

% \usepackage{scalerel}
% \newcommand{\omg}{\scalerel*{\includegraphics[height=1.5em]{figuras/omg-sorprendido-meme.png}}{X}}

\begin{multicols}{2}

% \begin{abstract}
%
% Os buracos negros son uns dos obxectos máis fascinantes do cosmos, xeralmente
% asociados á Teoría da Relatividade Xeral de Einstein. Porén, máis dun século
% antes de Einstein, dous científicos do século XVIII: John Michell e Pierre
% Simon Laplace, teorizaron de maneira independente sobre a existencia de
% estrelas tan masivas que nin sequera a luz podía escapar delas. Neste artigo
% exploramos como, empregando unicamente a física newtoniana e o concepto
% corpuscular da luz, ambos chegaron a conclusións sorprendentemente similares á
% nosa concepción moderna dos buracos negros.
%
% \end{abstract}

\subsection*{Introdución}

Os buracos negros son, sen ningunha dúbida, das cousas máis sorprendentes que
se descubriron en física o século pasado. A día de hoxe seguen sendo un enigma
e unha enorme fronte aberta en investigación. Teorizar con precisión os buracos
negros non tería sido posible sen a teoría da Relatividade Xeral de Einstein
nin as investigacións de Thorne, Wheeler e Hawking, entre moitos outros. Con
todo, como veremos a continuación, a historia sorpréndenos cun caso no que a
curiosidade humana permite idear conceptos tan adiantados ao seu tempo, que
requiriron dunha matemática tan avanzada para a súa resolución, que incluso os
propios matemáticos crían que se terían librado de nós, os físicos. Esta é a
historia de como John Michell e Pierre Simon Laplace formularon os primeiros
conceptos sobre buracos negros de maneira independente, máis dun século antes
da Relatividade Xeral.

\subsection*{Contexto histórico}

Ambas as formulacións teñen lugar nas décadas dos 80 e 90 do século XVIII. A
esas alturas da historia, Europa sufría os radicais cambios causados pola
Revolución Francesa e a caída dos reinos absolutistas. A Ilustración amosou os
seus froitos co triunfo do pensamento liberal. A física daquela época vivía
unha especie de período de transición entre dúas grandes revolucións
científicas: a revolución newtoniana do século XVII preparou o terreo para a
posterior revolución da electrodinámica e da estrutura da materia no século
XIX. Só 100 anos antes do inicio da nosa historia, Sir Isaac Newton publicaba
en 1687 \textit{Philosoph\ae~ Naturalis Principia Mathematica}, asentando as
bases da física moderna coas súas leis do movemento e da gravitación universal.
Dúas consecuencias deste desenvolvemento foron os protagonistas da nosa
historia: Michell e Laplace. Ambos chegaron a unha conclusión común:
\emph{existen estrelas tan masivas que fan curvar tanto a luz, que nin sequera
esta pode escapar da atracción gravitacional}~\footnote{Algo curioso desta
historia é que non hai un único descubridor deste fenómeno, xa que, ao existir
conflito entre Inglaterra e Francia, Laplace non soubo do artigo de Michell
cando saíu.} A continuación estudaremos como chegou cada un a esta conclusión.

\subsection*{Michell}

John Michell (Eakring Nottinghamshire, 1724) foi un matemático e reverendo
inglés. Destacou principalmente nos campos da xeoloxía e da astronomía. É
coñecido, ademais, como o pai da sismoloxía moderna, xa que foi o primeiro en
rexistrar os terremotos para a súa análise, ademais de inferir que se poderían
propagar por ondas sísmicas. Na astronomía, John Michell realizou un importante
traballo sobre estrelas binarias. Aplicando teoría da probabilidade a
distribucións estelares~\footnote{Sendo así o primeiro astrónomo estatístico.},
demostrou que a maioría dos sistemas binarios que se ven no espazo son como o
sistema \textit{Tatooine} de \textit{Star Wars} (estrelas que orbitan arredor
de outras) e non unha casualidade de perspectiva. Outra grande achega foi a
exitosa predición dos buracos negros, que é o obxecto deste artigo. Esta obra
publícase en 1784, baixo o nome \textit{<<excesivamente longo e
laborioso>>}.\cite{mongomeri}

\Cita{
    ``On the Means of Discovering the Distance, Magnitude {\&}c. of the Fixed
    Stars, in Consequence of the Diminution of the Velocity of Their Light, in
    Case Such a Diminution Should be Found to Take Place in any of Them, and
    Such Other Data Should be Procured from Observations, as Would be Further
    Necessary for That Purpose.''
}{Colin Montgomery}

Despois deste título practicamente non fai falta aportar moito máis, xa que
está case todo aí. Michell non só predixo a existencia dos buracos negros,
senón que ademais deu a ferramenta para detectalos. Ferramenta coa que se
detectou Cygnus X-1, o primeiro buraco negro descuberto e un temazo da banda
Rush. No artigo orixinal, Michell enuncia:

\Cita{
    ``...since their light could not arrive at us; [...], we could have no
    information from light; yet, if any other luminous bodies should happen to
    revolve about them we might still perhaps from the motions of these revolving
    bodies infer the existence of the central ones with some degree of
    probability...''
}{John Michell}

Agora ben, Michell non só elucubra que poderían existir buracos negros e como
detectalos. Senón que, empregando o valor estimado da velocidade da luz que se
coñecía entón, realiza unha análise das dimensións que debería ter unha estrela
para que se dean as condicións de \textit{estrela escura}. Michell empregou un
método xeométrico para realizar o estudo. Como se me fai incomprensible este
método, xa que dista radicalmente de calquera metodoloxía estudada en calquera
curso de física, vou explicar o procedemento en termos actuais que se nomea en
\cite[pág. 92]{mongomeri}\footnote{Igualmente animo o lector a botarlle un ollo
ao artigo orixinal~\cite{paper_largo_michell}, xa que ten certo interese
histórico ver como se realizaba a física antigamente. Tamén vén explicado
en~\cite[pág. 92]{mongomeri}.}.

\begin{figure}[H]
    \includegraphics[width=\linewidth]{revistas/004/imaxes/diagramita_michell.png}
    \caption{
        Diagrama orixinal empregado en \cite{paper_largo_michell} para a
        análise xeométrica. A figura \textit{per se} non se atopa no artigo
        arquivado pola Royal Society, pero atópase en \cite[Véxase nota número
        3]{mongomeri}.
    }
\end{figure}

\subsection*{A análise xeométrica de Michell (en termos modernos)}

Supoñamos dous corpos: un grande atractor de masa $M$ e raio $R$, e unha
partícula de masa $m$. O grande atractor (ou corpo central) e a partícula
atráense mutuamente mediante unha forza $F = G\frac{Mm}{r^2}$, onde podemos
substituír a densidade do corpo central, $\rho$
%
\begin{equation}
  F = K\frac{\rho m}{r^2}R^3.
\end{equation}

Onde $K = \frac{4\pi}{3}$. Podemos deducir a velocidade da masa $m$, que
corresponderá con
%
\begin{equation}
    v^2 = \frac{2K\rho R^3}{r} \quad\Rightarrow\quad v^2\propto
    \rho|_{r = \text{cte.}}
\end{equation}

Por tanto, imos fixar $r$ na superficie do corpo central, $r=R$,
%
\begin{equation}
  v^2 (r=R) = 2K\rho R^2,
\end{equation}
%
de onde deducimos unha curiosa propiedade: ao ser $2K$ unha constante, a
velocidade ao cadrado da partícula conservarase mentres non varíe o produto
$\rho R^2$. En termos matemáticos...
%
\begin{equation}
    \frac{v^2}{\rho R^2} = 2K = \text{const.}
\end{equation}

Esta $v$ será a velocidade orbital da partícula na superficie do corpo central.
Agora metámonos coa velocidade de escape, definida por Michell como
\textit{``the same veolcity a body would have if it fell from infinity to the
surface''}. Imos calculala mediante xogadas arcaicas que tivo que empregar
Michell na súa época. Para iso, imos meternos no sistema Terra-Sol.

A velocidade orbital da partícula na superficie do Sol será 14,65 veces a
velocidade orbital da Terra arredor do Sol
%
\begin{equation}
    v_{orbit}(r=R_S) = 14,65 v_{orbit}(r = r_{Terra}=1\,\text{UA}).
\end{equation}

Agora, a traxectoria dunha partícula que cae desde o infinito (i.e., enerxía
nula) será parabólica. Por tanto, a velocidade de escape na superficie do Sol
será
%
\begin{align}
    v_{parabol}(r=R_S) =& \sqrt{2} v_{orbit}(r=R_S) =\\
    =& 20,72 v_{orbit}(r = r_{Terra})
\end{align}
%
e, por tanto, $v_{parabol}(r=R_S) = v_{escape}(\text{Sol})$. Por que Michell
relaciona a velocidade de escape na superficie do Sol? Moi sinxelo, porque na
época a velocidade da luz coñecíase como 10310 veces a velocidade orbital da
Terra ao redor do Sol~\cite{luz_bradley} (Oh, sorpresa, a partícula de masa
% $m$ será un fotón {\large \omg}!). Deste xeito, podemos realizar unha
$m$ será un fotón {\large OMG}!). Deste xeito, podemos realizar unha
comparación directa entre a velocidade de escape e a velocidade da luz.
%
\begin{equation}
    \frac{c}{v_{escape}(\text{Sol})} = \frac{10310 v_{orbit}(r =
    r_{Terra})}{20,72 v_{orbit}(r = r_{Terra})} = 497.
\end{equation}

Por tanto, unha estrela será escura se a súa velocidade de escape é, cando
menos, 497 veces a velocidade de escape do Sol. Poderiamos dicir entón:
\emph{Se existise unha estrela coa mesma densidade que o Sol, pero cuxo raio
fose aproximadamente 500 veces maior, a luz non podería escapar dela e sería
invisible para nós.} En palabras de Michell:

\Cita{
    ``...if the semi-diameter of a sphere of the same density with the sun were
    to exceed that of the sun in the proportion of 500 to 1, a body falling
    from an infinite height towards it, would have acquired at its surface a
    greater velocity than that of light, and consequently, supposing light to
    be attracted by the same force in froportion to its vis inertiae, with
    other bodies, all light emitted from such a body would be made to return
    towards it, by its own proper gravity''
}{Michell \cite[Proposición 16, página 42]{paper_largo_michell}}

\subsection*{Laplace}

Pasamos agora a Laplace. Pierre Simon Laplace (Beaumont-en-Auge, 1749) foi un
matemático, astrónomo e físico francés. Foi un dos científicos máis influentes
da súa época, chegando a ser coñecido como o \textit{Newton francés}. A súa
obra máis coñecida é \textit{Mécanique Céleste} (1799–1825), na que reformula a
mecánica newtoniana para aplicala a sistemas astronómicos complexos. Nesta
obra, Laplace desenvolve a teoría nebular do sistema solar, que postula que o
sistema solar se formou a partir dunha nube de gas e po en rotación. Esta
teoría foi un precursor importante das teorías modernas sobre a formación de
sistemas planetarios. Laplace fixo todo iso e máis cousas que se ven na
carreira, polo que non as comentarei. Centrarémonos agora na súa predición dos
buracos negros.

Ao contrario de Michell, Laplace non expón todo nun só artigo, senón que, nun
primeiro momento, presenta a suposición de que poden existir este tipo de
corpos celestes e, anos máis tarde, divulga unha ``proba matemática''. O
contexto no que publica isto é, nada máis nin nada menos, o da Revolución
Francesa. Desde 1773, Laplace atopábase na Academia francesa, onde anos antes
aprendera física e matemáticas da man de D'Alembert.

En 1794 cae Robespierre e con el, a súa etapa do Terror. É entón cando o
sistema educativo francés experimenta un enorme cambio que conduce á
institucionalización da sociedade francesa moderna. Neste contexto fúndanse
diversas institucións científicas, incluíndo o \textit{Institut de France},
onde un ano despois Laplace sería elixido presidente. É por iso que cesa as
súas actividades como profesor na Academia e dedícase á redacción da súa obra
\textit{Exposition du Système du Monde}, publicada en dous volumes en 1796.

No sexto capítulo do quinto libro desta obra~\cite[página
348]{exposition_laplace} Laplace arrisca a enunciar cualitativamente a
existencia de estrelas tan masivas, que nin sequera a luz podería escapar
delas. Laplace escribe:

\Cita{
    ``Unha estrela luminosa da mesma densidade que a Terra, e cuxo diámetro
    fose duascentas cincuenta veces maior que o do Sol, non deixaría que ningún
    dos seus raios nos alcanzase en virtude da súa atracción; é, pois, posible
    que os maiores corpos luminosos do Universo sexan, por ese mesmo feito,
    invisibles.''
}{Laplace}

Ante a potente declaración de Laplace, un lector escéptico preguntouse: Como é
iso posible? Así, Franz Xaver von Zach escribiulle a Laplace unha carta na que
lle esixía unha demostración matemática da súa afirmación. Laplace responde e
en 1799 publícaa na revista alemá \textit{Allgemeine Geographische
Ephemeriden}, editada polo mesmo von Zach~\cite{laplace_original}. Tamén
existe unha tradución ao inglés desta demostración en~\cite{hawking_ellis}.

\subsection*{A ``proba matemática'' de Laplace}

Supoñamos unha estrela de masa $M$ e raio $R$, e un fotón de masa\footnote{Let
me cook} $m$, a unha distancia $r$ do centro da estrela.

Agora empregamos a conservación da enerxía en $r=R$ e nun $r$ arbitrario, tendo
en conta que a velocidade en $R$ sexa $c$, a velocidade da luz no baleiro. Por
tanto,
%
\begin{equation}
    \frac{mc^2}{2} - G\frac{Mm}{R} = \frac{mv^2}{2} - G\frac{Mm}{r}.
\end{equation}

Obviamente elimínanse as masas $m$. Agora particularizamos o cálculo no límite
$r\to\infty$, onde a velocidade do fotón é nula. É dicir, queremos que o noso
fotón chegue ao infinito e alí quede quieto. Esa é a nosa condición de contorno
no problema. Deste xeito
%
\begin{equation}
  \frac{c^2}{2} - G\frac{M}{R} = 0 \quad\Rightarrow\quad c = \sqrt{\frac{2GM}{R}},
\end{equation}
%
que é a moi coñecida formuliña da velocidade de escape. A clave do negocio
agora é o xenio de Laplace ao substituír a masa $M$ da estrela pola súa
densidade $\rho$ e o seu volume
%
\begin{equation}
  c = R\sqrt{\frac{8\pi}{3}\rho G},
\end{equation}
%
onde empregamos que $M = \frac{4}{3}\pi R^3 \rho$. Imos agora diverxer un pouco
do método orixinalmente empregado por Laplace. Este non desenvolve, como
Michell, o proceso matemático no que chega á conclusión das proporcións
necesarias para que unha estrela sexa escura. En lugar diso, Laplace, a partir
das proporcións enunciadas en~\cite{exposition_laplace}, amosa que a súa idea
ten sentido extraendo dela un valor coñecido. Neste caso, Laplace substituiría
os valores pertinentes para despexar da ecuación o valor da densidade
terrestre. Deste xeito asegurábase de que o lector comprendese que a súa
proposición tiña sentido. O problema dese método é que é moi aburrido. Do mesmo
xeito que con Michell, non concibo como comprender a súa metodoloxía orixinal.
Por tanto, farei unha exposición análoga, só que en vez de despexar a densidade
terrestre, despexarei a velocidade da luz empregando a ecuación anterior e os
valores propostos por Laplace en~\cite{exposition_laplace}.

Substituímos agora o raio da estrela por $R = 250 R_S$, onde $R_S = 6.9634\cdot
10^8$ m, e a densidade por $\rho = \rho_T = 5514$ kg/m$^3$. Por tanto
%
\begin{equation}
    c_{Laplace} = 250 R_S\sqrt{\frac{8\pi}{3}\rho_T G} = 3.057\cdot
    10^8 \text{ m/s}
\end{equation}

O cal ten un erro porcentual respecto do valor actual do 1,96\%. Nada mal.
Deste xeito demostramos, de maneira indirecta mediante a confirmación do valor
dunha magnitude coñecida, que a proposición de Laplace ten sentido.

\subsection*{Discusión}

Michell e Laplace foron dous xenios que conseguiron chegar a un resultado moi
adiantado ao seu tempo. Non é tarefa sinxela imaxinar como sería un obxecto tan
estraño como un buraco negro naquela época, e moito menos facelo por primeira
vez. Hai que comprender que, naquela altura, a mecánica newtoniana tiña os
mesmos anos que ten hoxe a mecánica cuántica, e aínda se debatía a natureza
corpuscular ou ondulatoria da luz. Ademais, se Young tivese feito o seu
experimento da dobre fenda antes de 1784, é probable que a ningún destes dous
se lle ocorrera a idea de que unha \textit{partícula} de luz fose atraída por
un corpo o suficientemente masivo.

O obxectivo deste artigo non só é aprender un pouco sobre a historia da física,
senón tamén intentar imaxinar o que pasaba polas mentes destes dous xenios.
Saír un día a deitarse a mirar as estrelas e visualizar o cosmos. Mirar a
constelación do Cisne e saber que, nas profundidades, se agocha unha forza
misteriosa e invisible: o buraco negro de Cygnus X-1. Por suposto, eles
descoñecían este astro escuro, pero invita a facerse unha idea do que tiña cada
un nas súas mentes.

Gustaríame rematar, unha vez máis, coa \textit{Exposition du Système du Monde}
de Laplace. Un pouco máis abaixo do enunciado das proporcións da estrela
escura, Laplace escribe:

\Cita{
    ``...finalmente, os movementos específicos de todos estes grandes
    corpos que, obedecendo a súa atracción mutua e, probablemente, a impulsos
    primitivos, describen inmensos orbes; tales serían, en relación coas
    estrelas, os principais obxectos da astronomía futura.''
}{Laplace}

\nocite{luz_bradley}
\nocite{laplace_original}
\nocite{star_wars_iv}
\nocite{cygnus_x1}

\printbibliography

\end{multicols}
